% =====================================================================
%  Основное содержимое отчёта.
% =====================================================================

\labpart{Описание установленного программного обеспечения}

\texttt{OSSEC} -- система обнаружения вторжений на основе хоста (HIDS),
устанавливаемая непосредственно на защищаемый сервер, а не работающая
на отдельном сетевом сенсоре (в отличие, например, от \texttt{Snort}
или \texttt{Suricata}). Задачи, которые она решает:

\begin{itemize}
  \item \textbf{Анализ журналов (log analysis)} -- построчное чтение
  указанных лог-файлов, разбор записи декодерами (для распознанных
  форматов, включая формат \texttt{sshd}) и сопоставление с базой
  правил обнаружения; при совпадении формируется тревога с уровнем
  критичности.
  \item \textbf{Контроль целостности файлов (syscheck/FIM)} --
  периодическая проверка контрольных сумм важных файлов системы на
  предмет несанкционированного изменения.
  \item \textbf{Обнаружение руткитов (rootcheck)} -- проверка признаков
  скомпрометированной системы (аномальные записи в системных
  структурах, известные сигнатуры).
  \item \textbf{Активный ответ (active response)} -- при тревоге выше
  заданного порога критичности автоматически выполняется команда
  реагирования: в данном случае \texttt{firewall-drop} (добавление
  правила \texttt{iptables DROP} для IP-источника) и \texttt{host-deny}
  (запись в \texttt{/etc/hosts.deny}), на 600 секунд. Именно эта
  возможность превращает систему из чисто наблюдающей (IDS) в
  способную самостоятельно вмешиваться (IPS).
\end{itemize}

\textbf{Принцип работы.} OSSEC установлен в режиме \texttt{local}
(сервер и агент совмещены на одной машине, что уместно для единственного
защищаемого хоста) из исходного кода (готового пакета для Ubuntu~24.04
нет). Отдельные процессы-демоны выполняют разные роли:
\texttt{ossec-logcollector} читает указанные файлы и передаёт новые
строки \texttt{ossec-analysisd}, который применяет декодеры и правила;
при срабатывании правила уровня~6 и выше \texttt{ossec-execd}
запускает соответствующий скрипт активного ответа. Отдельно настроен
список исключений (\texttt{allow\_list}) -- IP-адреса и подсети,
которые не блокируются активным ответом ни при каких обстоятельствах
(публичный адрес самого сервера и внутренняя подсеть, через которую
администратор подключается по VPN) -- без этого существует риск, что
система заблокирует легитимный, но интенсивный доступ, включая доступ
самого администратора.

\labpart{Анализируемые подсистемы и файлы}

Анализируемые подсистемы -- те же, что и в лабораторной работе №4.1:
служба SSH (\texttt{sshd}) и веб-сервер \texttt{Caddy}, -- с добавлением
самого OSSEC как третьей подсистемы, чьи собственные журналы теперь
тоже являются предметом анализа.

\begin{itemize}
  \item \textbf{\texttt{/var/log/auth.log}} и
  \textbf{\texttt{/var/log/caddy/access.log}} -- те же файлы, что
  описаны в отчёте по лабораторной работе №4.1; теперь они читаются не
  только вручную, но и потоково самим OSSEC.
  \item \textbf{\texttt{/var/ossec/logs/alerts/alerts.log}} -- журнал
  всех тревог OSSEC с указанием сработавшего правила, его уровня
  критичности, исходного события и, если оно было определено,
  IP-адреса источника.
  \item \textbf{\texttt{/var/ossec/logs/active-responses.log}} --
  журнал именно исполненных действий активного ответа: какой IP,
  когда и на какой срок был заблокирован. Это основной источник данных
  для демонстрации перехода от обнаружения к предотвращению.
\end{itemize}

\labpart{Анализ запросов из сети Internet до и после настройки}

\begin{quote}
\textit{Раздел будет заполнен после нескольких дней работы сервера с
включённым OSSEC. Данные лабораторной работы №4.1 (период без
активного ответа) послужат базовым уровнем «до»; здесь будет
приведено прямое сравнение с периодом «после»: изменение числа
успешных подключений с уже заблокированных источников, число и состав
реальных срабатываний активного ответа (\texttt{active-responses.log}),
а также то, повлияла ли работа OSSEC на журнал веб-сервера (учитывая,
что активный ответ в текущей конфигурации настроен по правилам,
специфичным для SSH, а не для HTTP-трафика).}
\end{quote}

\labpart{Сравнение с аналогами}

\begin{table}[H]
\centering
\small
\begin{tabular}{|p{1.8cm}|p{2.9cm}|p{3.2cm}|p{4.6cm}|}
\hline
\textbf{Система} & \textbf{Тип} & \textbf{Источник данных} & \textbf{Особенность} \\
\hline
OSSEC & HIDS (+active response) & Логи файлов на хосте & Лёгкая установка на один сервер, встроенный IPS-режим через скрипты \\
\hline
Snort & NIDS/NIPS & Сетевые пакеты (libpcap) & Классический сигнатурный анализ трафика, требует описания правил под протокол \\
\hline
Suricata & NIDS/NIPS & Сетевые пакеты, многопоточность & Совместим с правилами Snort, дополнительно разбирает прикладные протоколы (HTTP, TLS и т.\,д.) \\
\hline
Tripwire & HIDS (только FIM) & Контрольные суммы файлов & Узкоспециализирован на контроле целостности, не анализирует сетевой трафик и логи в реальном времени \\
\hline
Zeek (Bro) & NIDS (анализ, не блокировка) & Сетевые пакеты & Ориентирован на подробное протоколирование и последующий анализ, а не на немедленную блокировку \\
\hline
\end{tabular}
\caption{OSSEC и альтернативные варианты задания}
\end{table}

Для единственного сервера с ограниченным набором подсистем (SSH и один
веб-сервер) выбор HIDS-решения, читающего уже существующие логи, вместо
NIDS-решения, анализирующего сырой сетевой трафик, обоснован
меньшими требованиями к настройке (не нужно описывать сигнатуры для
протоколов) и меньшей нагрузкой на единственный процессор сервера.

\labpart{Демонстрация работы установленного программного обеспечения}

\begin{quote}
\textit{Раздел будет дополнен снимками экрана/выдержками из
\texttt{active-responses.log} после того, как накопится статистика
реальных срабатываний.}
\end{quote}
